% ===================================================================
\section{Conclusion}\label{sec:conclusion}
% ===================================================================
We estimated the Brent-price effect of the 2026 Strait of Hormuz disruption by building a
counterfactual with an ensemble synthetic control and benchmarking the design against the
documented Russia 2022 shock. Over the window to late June 2026 the ensemble-median premium
is roughly \SI{53}{\percent} of Brent's no-event level (IQR 48 to 53\%), large and positive
under every estimator. 

Brent ranks first among all units in the in-space placebo at the attainable floor, the
leave-one-out check stays tight, and the contamination-bias sign points conservative, so the
estimate clears the validation battery. Outside that battery, the synthetic beats the
donor-free naive baselines, and on the benchmark event the same machinery recovers a magnitude
consistent with the documented \$95 to \$130 move and within \$3.6 of the EIA's pre-invasion
structural forecast, independent reassurance that the method tracks a real path.

The substantive result is a data-driven Brent premium for the Hormuz disruption that
complements the structural scenario estimates produced by energy institutions, and the design
that recovers it is itself a result, a transparent template for validating cross-asset
synthetic control under a single-unit geopolitical treatment, with an identification argument a
reader can audit donor by donor. The structural approach and ours answer different questions,
the structural models asking what price an assumed barrel shortfall implies and the synthetic
control asking what price gap the observed co-movement of global markets implies. Because the
two read the event from different directions, they work best together, and our estimate is a
complement to scenario analysis rather than a substitute for it.

With a single treated unit and a short post-period, permutation inference has little power to
draw on, and the Hormuz placebo reaches 0.050, the
best value attainable at this sample size rather than a conventional rejection, so we read it
as the strongest available signal and not as a small $p$-value in the usual sense. The estimand
is also a reduced-form total effect, capturing the full price gap but not separating it into a
physical-supply, a risk-premium, and an expectations component that some policy uses would want
to distinguish. Being recent, the Hormuz event has no documented post-event counterfactual of
its own, so its validation is internal and rests on the Russia case as a calibration by analogy
rather than on an external ground truth. The factor structure, too, is regime-bound, as the
weak cross-event transfer shows directly, so the estimate holds for the regime of this event
and should not be read as a general chokepoint premium that would recur unchanged in a
different macro environment.

Beneath all of these, donor cleanliness is only partly testable. The no-interference
assumption concerns an unobserved counterfactual donor path, and the contamination-sign
diagnostic is naive and directional, so it bounds the likely direction of any residual bias
without proving that none exists. None of this overturns the central finding. The observed
co-movement of global markets with Brent implies that the 2026 Hormuz disruption added a large
and robust premium to crude, on the order of half its no-event price, and the estimate is
meaningful and usable for the decisions that depend on it.


